\documentclass[12pt]{article}
\usepackage[utf8]{inputenc}
\usepackage[T1]{fontenc}
\usepackage{hyperref}
\title{Hod\"i Language Resources for K-12 Curriculum Development}
\author{Compiled by Claude Sonnet 5.5 (claude-sonnet-5-5) for Andr\'as Kornai}
\date{}
\begin{document}
\maketitle

\section{Introduction}
Hod\"i (also J\"oti, Yuwana, Waruwaru, Chikano) is a language of the Venezuelan Amazon with no firmly established relatives. Some proposals link it to Piaroa--Saliban, but others treat shared vocabulary as loans. Speakers number about 640 (2007). The 2011 census counted 982 Hod\"i people. They call themselves Jojod\"o, and the language Jojod\"o tjiwene (as given in Wikipedia mirrors). Only one search was run for this language, and it returned mostly encyclopedic pages. No educational materials were found.

\section{Primary Educational Resources}
No primers, textbooks or school programmes were found.

\section{Linguistic and Reference Materials}
\begin{itemize}
\item Wikipedia article: \url{https://en.wikipedia.org/wiki/Hod%C3%AF_language}
\item ISO 639:yau Wikipedia page: \url{https://wikipedia.com/wiki/ISO_639:yau}
\item Endangered Languages Catalogue entry: \url{https://endangeredlanguages.com/lang/2902}
\item Verbix endangered languages page on J\"oti: \url{https://docs.verbix.com/EndangeredLanguages/Joti}
\item EcuRed article on the Hoti people (in Spanish): \url{https://www.ecured.cu/Hoti}
\end{itemize}

\section{Religious and Community Resources}
No scripture translations or community organisations were found in the search results.

\section{Digital and Audio Resources}
No audio, video or software resources were found.

\section{Curriculum Development Context}
The language is stable and used by the whole community, but nothing found indicates a written standard, teaching materials or schooling in Hod\"i. Any curriculum would have to begin with documentation and orthography work, led by the Hod\"i community and in line with Venezuela's intercultural bilingual education policy, which I did not verify here. Further searches for linguistic descriptions and recordings are advisable before planning.

\end{document}